\section{Conclusion}

A furlong's flight from the abbey tower is physically possible for a cloth wing, but only
in a strong southwest wind, with a wing near the better end of what the materials allow,
and on a track that follows the slope.
A flight straight out over the valley does not reach that distance in any sampled case.
Without a tail the craft is unstable in pitch, with a disturbance doubling in about the
time a man takes to react, so a flight of some tens of seconds followed by a loss of control
in a gust fits the models; a tail removes the instability.
The broken legs require a fast, steep arrival, which a loss of control provides.
Eilmer's own explanation of his fall is consistent with the physics, though the physics
cannot show that it was the only cause.

\section*{Data and code availability}

\ifdefined\anon
The analysis code, data manifests, and figure scripts reproduce every number and figure in
this paper.
They are openly available in a public repository archived with a DOI; the links are
withheld for double-anonymous review and have been supplied to the editors.
\else
The analysis code, data manifests, and figure scripts reproduce every number and figure in
this paper.
They are openly available, and archived with a DOI, at
\begin{center}
\url{https://github.com/silverdavi/eilmer-flight-feasibility}\\
\url{https://doi.org/10.5281/zenodo.20706321}
\end{center}

\section*{Acknowledgements}

The author thanks the Athelstan Museum, Malmesbury, for its local record of Eilmer's flight,
and Dr Duncan Pepper for the photograph in \cref{fig:photo}, published under a CC BY-SA 2.0
licence.
This research received no specific funding.
\fi
